\documentclass[10pt,a4paper]{amsart}
\usepackage[margin=19mm]{geometry}
\usepackage{amsmath,amssymb,mathtools}
\usepackage[scr=rsfs,cal=euler]{mathalfa}
\usepackage{xfrac}
\usepackage[unicode,hidelinks,bookmarks=false]{hyperref}
\newcommand{\ie}{i.e.\ }
\newcommand{\cf}{cf.\ }
\newcommand{\resp}{respectively\ }
\newcommand{\Subsection}{\subsection}
\providecommand{\nobreakdash}{\nobreak\hbox{-}}
\setlength{\emergencystretch}{2em}
\multlinegap0pt
\usepackage{bm}
\usepackage{tikz}
\usepackage{amssymb}
\usepackage{mathtools}
\usepackage{algorithm}\usepackage{algpseudocode}
\newtheorem{lem}{Lemma}
\newtheorem{prop}{Proposition}
\newcommand{\B}{\mathcal{B}}
\newcommand{\C}{\mathcal{C}}
\newcommand{\G}{\Gamma}
\newcommand{\Wh}{\mathrm{Wh}}
\newcommand{\lk}{\mathrm{lk}}
\newcommand{\s}{\mathcal{S}}
\newcommand{\st}{\mathrm{st}}
\title{Cohomology of pure symmetric automorphisms of right-angled Artin groups}
\author{Peio Ardaiz-Gale}
\date{Source: arXiv:2604.13749v3 (CC BY 4.0)}
\begin{document}
\maketitle

\noindent\emph{Source and licence.} Cohomology of pure symmetric automorphisms of right-angled Artin groups. Version arXiv:2604.13749v3 (CC BY 4.0), distributed by arXiv under the Creative Commons Attribution 4.0 International licence, http://creativecommons.org/licenses/by/4.0/, as recorded in the arXiv licence metadata for this exact version and shown on the abstract page https://arxiv.org/abs/2604.13749v3. Full licence notice: \texttt{licenses/arxiv-2604.13749-CC-BY-4.0.txt}.\par\smallskip

\noindent\textit{Editorial scope.} Counted unit: the article's prop prop:baseIsomor (source0.tex lines 1273--1275). The article's complete original proof of it is delivered with it (source0.tex lines 1276--1390, one \texttt{proof} environment of 115 lines). No mathematics is written, repaired or re-derived: the statement and the proof are the article's own lines, in the article's own notation, and no retained line is substituted or re-flowed.  Retained uncounted as the source-local context the proof needs: lines 1257--1259; lines 1266--1268.  Nothing was omitted from any retained span, so the package carries no omission record.  No retained line contains a citation, so the wrapper carries no bibliography and no bibliography entry.  No retained line contains a figure environment, an image-inclusion command or drawing code, so no image is delivered and the package declares no asset dependency.  Editorial formatting only, in addition: an AMS \texttt{amsart} preamble that restates the article's own declarations and macros used by the retained lines, namely the environments \texttt{lem}, \texttt{prop} on separate counters, together with the standard packages those lines need.  The retained environments print in this document as Lemma 1, Proposition 1 in order of appearance, whereas the article prints the same items with the numbering of its own full document.  The complete native source of the article as distributed in the arXiv e-print for this exact version is bundled unchanged as \texttt{sources/198495/source0.tex}.
\begin{lem}\label{lem:minimal-component}
    Let $v_i\notin\st(v_j)$ and let $\min_i$ denote the minimal element in $\G-\st(v_i)$. Let $\mathcal{S}$ denote a subordinate component.  If $\min_i\neq\min_j$ then either $\min_i\in D^j$, or $\min_j\in D^i$ or both. If  $\min_i=\min_j$ then either $\min_i\in D^j$ and $\min_j\in \mathcal{S}\subset D^j$,  or $\min_j\in D^i$ and $\min_i\in \mathcal{S}\subset D^i$,  or $\min_i$ and $\min_j$ are in the same shared component $\C$.
\end{lem}

\begin{lem}\label{lem:rank2Partitions}
    Let $\tau$ be a rank 2 essential vertex type in $\Wh_\G$ with two non trivial based partitions $\tau_i=\{\{v_i\}, P_1, P_2\}$ and $\tau_j=\{\{v_j\}, Q_1, Q_2\}$. Assume that $\min_i\in P_1$ and $\min_j\in Q_1$. Then either $P_2$ and $Q_2$ both consist of a unique connected component, or $P_2=\C$ is a shared component and $Q_2=D^i\cup\C$ (or the other way around).
\end{lem}

\begin{prop}
    There exists an surjective map   $\varphi: \B_1\rightarrow\B_2$. \label{prop:baseIsomor}
\end{prop}

\begin{proof}
We will start by  constructing $\varphi$ explicitly for each of the types of elements of $\B_1$ considering several possible cases. During the construction, we will denote by boldface  $\boldsymbol{A}$ the petal $A$ containing the minimal element $\min_i$ in the based partition $\tau_i$. 

    \noindent{\textbf{Type 1:}} Rank 2 vertex types $\tau$.

     {\textbf{1.1:}} $\tau$ has a single non-trivial based partition $\tau_i=\{\{v_i\},P_1,P_2,P_3\}$:  We may assume that  $\min_i\in P_1$. Then, since $\tau$ is essential, $P_2$ and $P_3$ cannot admit a split, so they must consist of a unique connected component, because otherwise they could be split, since all of the other based partitions are trivial. We set 
        $$
        \varphi(\tau)= \gamma^i_{P_2}\gamma^i_{P_3}\in\B_2.
        $$

     
     {\textbf{1.2:}} $\tau$ has two non-trivial based partitions $\tau_i=\{\{v_i\},P_1,P_2\}$ and $\tau_j=\{\{v_j\},Q_1,Q_2\}$:
     We  assume that that $\min_i\in P_1$ and $\min_j\in Q_1$. By Lemma \ref{lem:rank2Partitions}, we have two possible cases:
        
         {\textbf{1.2.1:}}  $v_i\notin \lk(v_j)$, $P_2=\C$ a shared component and $Q_2=D^i\cup\C$. We set
         $$\varphi(\tau)= \gamma^i_{P_2}\gamma^j_{D^i}\in\B_2$$

         {\textbf{1.2.2:}} In other case, since $P_2$ and $Q_2$ are formed by a unique connected component, we set:
        $$\varphi(\tau)= \gamma^i_{P_2}\gamma^j_{Q_2}\in\B_2.$$
        Note that $P_2$ and $Q_2$ must be different, because if $P_2=\C=Q_2$ there would be a crossing.

         
       % Let us check that the claim in 1.2.1. We know that $P_2$ and  $Q_2$ cannot be split. Since $P_2$ does not contain the dominant, then it must consist of a unique connected component.  On the other hand, $Q_2=\{D^i\cup \mathcal{C}_1\cup\cdots \cup \mathcal{C}_k\}$, where each $\mathcal{C}_i$ is shared, because subordinate components do not produce crossings and can always be split. But these $\mathcal{C}_1,...,\mathcal{C}_k$ must be in $P_2$, to impede the split of $Q_2$ being valid. As a consequence, $P_2=\C$ and $Q_2=D^i\cup \C$. Since $P_2$ consists of a unique connected component which is not $D^j$ we have that $\gamma^i_{P_2}\gamma^j\in \B_2$.

       % In  case 1.2.2, we have already seen that $P_2$ and $Q_2$ cannot contain  both the respective dominant components. So it remains to check that both $P_2$, $Q_2$ consist of a unique connected component and that $P_2\neq Q_2$, i.e., they are not shared and equal. We have three possibilities
          %  \begin{itemize}
           %     \item   $v_i\in \st(v_j)$. In this case there cannot be crossings between $\tau_i$ and $\tau_j$, so any split is allowed. As a consequence, since they are essential, both $P_2$ and $Q_2$ must be formed by a unique connected component; since they are adjacent there are no shared  components so $\gamma^i_{P_2}\gamma^j_{Q_2}\in \B_2$.

              %  \item  $v_i\notin \st(v_j)$ and both $P_1$ and $Q_1$ contain the respective dominant components.   Since $P_2$ and $Q_2$ cannot admit a split, both $P_2$ and $Q_2$  must be formed by a unique connected component; they can not be equal since we would have a crossing  so  $\gamma^i_{P_2}\gamma^j_{Q_2}\in \B_2$.

             %   \item  $v_i\notin \st(v_j)$ and $P_1$ contains the dominant component and $Q_2$ is just the dominant component. By the same argument as before, we know that $P_2$ must be formed by a unique connected component. $Q_2$ is also formed by just one connected component and it is not shared so $\gamma_ {P_2}^j\gamma_{Q_2}^i\in \B_2$.
           % \end{itemize}

            
            
         
 


\noindent %It is clear that the map $\varphi$ constructed so far is injective. Besides,
%Notice that for all of the elements of the form $\gamma^i_{A}\gamma^i_{B}$ or $\gamma^i_{A}\gamma^j_{B}$ that we have in the image of Type 1 elements the   $A,B$ never contain the respective minimal element. 


\noindent{\textbf{Type 2:}} Pairs $(\tau,v_j)$ for $\tau$ a rank 1 essential vertex and $v_j\in\G$. In $\tau$ every based partition is trivial except for  $\tau_i=\{\{v_i\},P_1,P_2\}$. We  assume that   $\min_i\in P_1$, so $P_2$ must consist of a unique connected component; in other case, it would admit a split. Again, we consider several subcases. 

{\textbf{2.1:}}  Assume $v_i\notin \st (v_j)$,  $\min_i\neq\min_j$ one of the $\min_i,\min_j$ is in a shared component and the other in a dominant component.

{\textbf{2.1.1:}}  If $\min_j\in\C$, a shared component, then, by Lemma \ref{lem:minimal-component},  $\min_i\in D^j$. For the pair $(\tau,v_j)$ with $\tau_i=\{\{v_i\},P_1,\C\}$, we set:
$$
\varphi(\tau,v_j)=\gamma^i_{\boldsymbol{D}^j}\gamma^j_{\boldsymbol{\C}}
$$

{\textbf{2.1.2:}} If $\min_i\in\C$, a shared component, then, by Lemma \ref{lem:minimal-component},  $\min_j\in D^i$. For the pair $(\tau,v_j)$ with $\tau_i=\{\{v_i\},P_1,D^j\}$, we set:
$$
\varphi(\tau,v_j)=\gamma^i_{D^j}\gamma^j_\C
$$

{\textbf{2.2:}} Assume $v_i\notin \st (v_j)$, $\min_j\in D^i$,   $\min_i\in\mathcal{S}$, a subordinate component, and the only non-trivial based partition in $\tau$ is $\tau_i=\{\{v_i\},P_1,D^j\}$. We set:
 $$
    \varphi(\tau,v_j)= \gamma^i_{\boldsymbol{\s}}\gamma^j_{\boldsymbol{D}^i}
    $$

    {\textbf{2.3:}} In any other case, for the pair $(\tau,v_j)$ with $\tau_i=\{\{v_i\},P_1,P_2\}$ such that  $\min_i\in P_1$ and  $\min_j\in B$, we set:
    $$
    \varphi(\tau,v_j)= \gamma^i_{P_2}\gamma^j_{\boldsymbol{B}}
    $$

    
\noindent{\textbf{Type 3:}} Edges $(v_i,v_j)\in E(\G)$. Consider the connected components $A$ of $\G-\st(v_i)$ and $B$ of $\G-\st(v_j)$  that contain the minimal elements and set
 $$
    \varphi(i,j)= \gamma^i_{\boldsymbol{A}}\gamma^j_{\boldsymbol{B}}
    $$



During the construction we have already checked that the elements chosen in the image of $\varphi$ lie in $\B_2$; equivalently,  that they are of the form $\gamma^i_{A}\gamma^j_{B}$ where  $A$ and $B$ consist of a unique  connected component in $\G-\st(v_i)$ and $\G-\st(v_j)$ respectively, $A\neq B$, and $A,B$ are not the respective dominant components at the same time. 

It remains to prove that the map is surjective. Let $\gamma^i_{A}\gamma^j_{B}\in\B_2$.

  %check that every $\gamma^i_{A}\gamma^j_{B}\in\B_2$ is the image of some element in $\B_1$.


\noindent{\textbf{Case I:}}  Assume $\min_i\notin A$ and $\min_j\in B$. Let $\tau$ be the rank 1 vertex type with only non trivial based partition  $\tau_i=\{\{v_i\}, A^C, A\}$, here $A^C$ denotes $(\G-\st(v_i))-A$. We have that $\tau$ is essential, since $A$ consists of a unique connected component. Hence, $(\tau,v_j)$ is of type 2.3 and $\varphi(\tau,v_j)=\gamma^i_{A}\gamma^j_{B}$.

\noindent{\textbf{Case II:}} Assume $\min_i\in A$ and $\min_j\in B$ (so $v_i\neq v_j$).
\begin{itemize}
    \item If $v_i\in\lk(v_j)$, then $(v_i,v_j)$ is of type 3 and $\varphi(v_i,v_j)=\gamma^i_{A}\gamma^j_{B}$.
    \item If $v_i\notin\lk(v_j)$, then $A$ and $B$ cannot be shared and equal, and they cannot be both dominant. By Lemma \ref{lem:minimal-component}, one must be dominant and the other is either shared or subordinate.

    (1) Assume $A=D^j$ dominant and $B$ shared. Let $\tau$ be the rank 1 vertex type with only non trivial based partition  $\tau_i=\{\{v_i\}, B^C, B\}$. We have that $\tau$ is essential, since $B$ consists of a unique connected component (note that $\min_i\notin B$). Hence, $(\tau,v_j)$ is of type 2.1.1 and $\varphi(\tau,v_j)=\gamma^i_{A}\gamma^j_{B}$.

    (2) Assume $B=D^i$ dominant and $A$ subordinate. Let $\tau$ be the rank 1 vertex type with only non trivial based partition  $\tau_i=\{\{v_i\}, (D^j)^C, D^j\}$. Again, $\tau$ is essential. Hence, $(\tau,v_j)$ is of type 2.2 and $\varphi(\tau,v_j)=\gamma^i_{A}\gamma^j_{B}$.
\end{itemize}

\noindent{\textbf{Case III:}} Assume $\min_i\notin A$ and $\min_j\notin B$.
\begin{itemize}
    \item If $v_i=v_j$, let  $\tau$ be the rank 2 essential vertex type with only non trivial based partition  $\tau_i=\{\{v_i\}, (A\cup B)^C, A, B\}$. Then, $\tau$ is of type 1.1 and $\varphi(\tau)=\gamma^i_{A}\gamma^i_{B}$.
    \item Assume that  $\tau_i=\{\{v_i\}, A^C, A\}$ and $\tau_j=\{\{v_j\}, B^C, B\}$ are compatible. Then the rank 2 vertex type $\tau$ with only non trivial based partitions  $\tau_i$ and $\tau_j$ is essential, so it is of type 1.2.2 and $\varphi(\tau)=\gamma^i_{A}\gamma^j_{B}$.
    \item In other case: the fact that $\tau_i$ and $\tau_j$ are not compatible implies that one of $A,B$  must be  dominant and the other shared. We assume that $A=D^j$, so $B$ is shared.

    (1) If  $\min_i\in B$, let $\tau$ be the rank 1 essential vertex type with only non trivial based partition  $\tau_i=\{\{v_i\}, A^C, A\}$. Then, $(\tau,v_j)$ is of type 2.1.2 and $\varphi(\tau,v_j)=\gamma^i_{A}\gamma^j_{B}$.

    

    (2) If  $\min_i\notin B$, let $\tau$ be the rank 1 vertex type with only non trivial based partitions $\tau_i=\{\{v_i\}, (A\cup B)^C, A\cup B\}$ and $\tau_j=\{\{v_j\}, B^C,  B\}$. The petal $A$  does not admit a split because a crossing would appear, so $\tau$ is essential. Hence, $\tau$ is of type 1.2.1 and $\varphi(\tau)=\gamma^j_{B}\gamma^i_{A}$.

\end{itemize}







\end{proof}

\end{document}
